\documentclass[12pt]{article}
\usepackage[utf8]{inputenc}
\usepackage[T1]{fontenc}
\usepackage[margin=1in]{geometry}
\usepackage{setspace}
\usepackage{xcolor}
\usepackage[colorlinks=true, urlcolor=blue!60!black, linkcolor=blue!60!black]{hyperref}
\setlength{\parindent}{2em}
\setlength{\parskip}{1.25em}
\renewcommand{\baselinestretch}{1.0}
\title{\textit{United States v. Butler}, 297 U.S. 1 (1936)}
\author{Case Brief}
\date{}

\begin{document}
\maketitle

\section{Case Metadata}
\begin{itemize}
  \item \textbf{Citation:} \textit{United States v. Butler}, 297 U.S. 1 (1936) (No.\ 401).
  \item \textbf{Argued / Decided:} Argued December 9--10, 1935; Decided January 6, 1936.
  \item \textbf{Vote:} 6--3 (affirming the First Circuit).
  \item \textbf{Majority Author:} Justice Owen J.\ Roberts (joined by Chief Justice Charles Evans Hughes and Justices Willis Van Devanter, James C.\ McReynolds, George Sutherland, and Pierce Butler).
  \item \textbf{Concurrences:} None.
  \item \textbf{Dissents:} Justice Harlan Fiske Stone (joined by Justices Louis D.\ Brandeis and Benjamin N.\ Cardozo).
  \item \textbf{Chapter / Topic:} Federal Legislative Power --- The Taxing and Spending Power (Art.\ I, \S\ 8, cl.\ 1) and Tenth Amendment Federalism Limits.
  \item \textbf{Primary Source URLs:} \href{https://supreme.justia.com/cases/federal/us/297/1/}{Justia U.S.\ Supreme Court Center (297 U.S.\ 1)}; \href{https://www.law.cornell.edu/supremecourt/text/297/1}{Cornell LII (297 U.S.\ 1)}.
\end{itemize}

\section{Facts and Procedural History}
During the Great Depression, collapsing farm prices created a severe disparity between the purchasing power of agricultural commodities and industrial goods. To restore farmers' purchasing power to its pre-World War I base-period level (August 1909 to July 1914 for cotton and most covered crops), Congress enacted Title I of the Agricultural Adjustment Act of May 12, 1933, ch.\ 25, 48 Stat.\ 31 (``AAA''). Section 8 of the AAA authorized the Secretary of Agriculture to enter into voluntary agreements with farmers to reduce acreage or crop production in exchange for ``rental or benefit payments.'' To finance these payments, \S\ 9(a) levied a ``processing tax'' on the first domestic processing of covered basic agricultural commodities (including cotton), set at a rate equal to the difference between the current average farm price and the statutory ``fair exchange value,'' while \S\ 16 imposed an equivalent ``floor tax'' on existing inventories of processed articles. Section 12(b) earmarked the proceeds of these taxes specifically for the Secretary of Agriculture to fund rental and benefit payments, market expansion, surplus removal, administrative costs, and tax refunds.

On July 14, 1933, the Secretary of Agriculture proclaimed that rental and benefit payments would be made for cotton starting August 1, 1933, and fixed the corresponding processing and floor tax rates. The United States subsequently presented a tax claim under \S\S\ 9 and 16 to William M.\ Butler and the other receivers of the insolvent Hoosac Mills Corporation, a Massachusetts cotton-milling company. The receivers recommended disallowing the claim on constitutional grounds. The U.S.\ District Court for the District of Massachusetts upheld the taxes and ordered them paid (\textit{Franklin Process Co.\ v.\ Hoosac Mills Corp.}, 8 F.\ Supp.\ 552 (D.\ Mass.\ 1934)). The U.S.\ Court of Appeals for the First Circuit reversed, holding the processing and floor taxes unconstitutional (\textit{Butler v.\ United States}, 78 F.2d 1 (1st Cir.\ 1935)). The Supreme Court granted certiorari (296 U.S.\ 561) and affirmed.

\section{Issue(s) Presented}
\begin{enumerate}
  \item Whether processors assessed earmarked processing and floor taxes under the Agricultural Adjustment Act of 1933 have standing to challenge the constitutionality of the spending program funded by those taxes, notwithstanding \textit{Massachusetts v.\ Mellon}, 262 U.S.\ 447 (1923).
  \item Whether the General Welfare Clause of Article I, \S\ 8, cl.\ 1 grants Congress an independent, substantive power to tax and spend for the national welfare beyond its other enumerated legislative powers (the Hamiltonian view), or instead confines spending to the other enumerated powers in Article I, \S\ 8 (the Madisonian view).
  \item Whether Congress may use its taxing and spending power under Article I, \S\ 8, cl.\ 1 to levy an earmarked processing tax and pay subsidies to farmers who contractually agree to curtail agricultural production, or whether such a scheme unconstitutionally invades powers reserved to the states under the Tenth Amendment.
\end{enumerate}

\section{Holding}
\begin{enumerate}
  \item \textbf{Standing:} Yes. Unlike a general federal taxpayer challenging an unearmarked appropriation from the general Treasury, a processor resisting an earmarked exaction that is an integral part of a single regulatory scheme has standing to challenge the validity of the entire plan.
  \item \textbf{Scope of the General Welfare Clause:} The Court adopted the \textbf{Hamiltonian (and Story) view}: Article I, \S\ 8, cl.\ 1 confers a substantive power to tax and appropriate for the ``general Welfare of the United States'' that is separate from, and not confined to, the other direct grants of legislative power in Article I, \S\ 8.
  \item \textbf{Constitutionality of the AAA:} Unconstitutional (6--3). Even under the broad Hamiltonian reading of the Spending Clause, the Agricultural Adjustment Act's processing tax and acreage-reduction payment scheme violated the Tenth Amendment because it was a pretextual plan to coerce and purchase compliance with federal regulation of local agricultural production---a subject reserved exclusively to the states.
\end{enumerate}

\section{Reasoning}
Writing for a six-Justice majority, Justice Roberts structured the opinion in three analytical steps: standing and the nature of the exaction, the constitutional scope of the Taxing and Spending Clause, and Tenth Amendment limits on regulatory spending.

\paragraph{1. Standing and the Non-Tax Character of the Exaction.}
The Government argued under \textit{Massachusetts v.\ Mellon}, 262 U.S.\ 447 (1923), that Hoosac Mills's receivers lacked standing to challenge how federal tax revenues were spent once collected. Justice Roberts distinguished \textit{Mellon} on the ground that the receivers were not general taxpayers objecting to a minute, indeterminable share of Treasury disbursements; rather, they were directly resisting an exaction that was ``an indispensable part in the plan of regulation'' (297 U.S.\ at 59). Because the tax automatically began and ended with benefit payments, was calibrated to parity prices, and was earmarked solely for crop control rather than general governmental support, it was not a true ``tax'' in the strict constitutional sense, which ``signifies an exaction for the support of Government'' and not ``the expropriation of money from one group for the benefit of another'' (id.\ at 61). Yet, echoing the \textit{Head Money Cases}, 112 U.S.\ 580 (1884), Roberts acknowledged that an earmarked exaction could still survive if it were an incident of a regulation within a valid enumerated federal power.

\paragraph{2. Judicial Review and the Madison vs.\ Hamilton General Welfare Debate.}
Before turning to Article I, \S\ 8, Justice Roberts articulated his famous formalist description of judicial review: a court neither approves nor condemns legislative policy, but has ``only one duty---to lay the article of the Constitution which is invoked beside the statute which is challenged and to decide whether the latter squares with the former'' (297 U.S.\ at 62). Because the Government did not attempt to defend the AAA under the Commerce Clause---agricultural production then being viewed as ``a purely local activity'' (id.\ at 64)---the statute stood or fell on Article I, \S\ 8, cl.\ 1, which empowers Congress ``to lay and collect Taxes, Duties, Imposts and Excises, to pay the Debts and provide for the common Defence and general Welfare of the United States.''

Roberts first confirmed that the phrase ``to provide for the general welfare'' is not an independent grant of general police power, but a qualification defining the purposes for which Congress may tax and (by necessary implication under Art.\ I, \S\ 9, cl.\ 7) appropriate funds. He then resolved the century-and-a-half-old Founding-era dispute between James Madison and Alexander Hamilton:
\begin{itemize}
  \item \textbf{The Madisonian View:} Madison contended that ``general welfare'' merely referred back to the specific powers subsequently enumerated in Article I, \S\ 8, confining federal taxation and spending to those enumerated fields. Roberts rejected this reading as rendering the General Welfare Clause ``mere tautology'' (id.\ at 65).
  \item \textbf{The Hamiltonian (and Story) View:} Hamilton---later endorsed by Justice Joseph Story in his \textit{Commentaries on the Constitution}---maintained that the clause confers a ``separate and distinct'' substantive power to tax and appropriate, limited only by the requirement that the expenditure serve the \textit{national} (as opposed to purely local) general welfare.
\end{itemize}
Siding with Hamilton and Story, the Court held that ``the power of Congress to authorize expenditure of public moneys for public purposes is not limited by the direct grants of legislative power found in the Constitution'' (id.\ at 66).

\paragraph{3. Tenth Amendment Barrier: Coercion and ``Purchased'' Regulation.}
In a dramatic doctrinal pivot, however, Justice Roberts declined to decide whether aid to agriculture promoted the ``general welfare'' (id.\ at 68). Instead, invoking the Tenth Amendment and Chief Justice Marshall's pretext warning in \textit{McCulloch v.\ Maryland}, 17 U.S.\ (4 Wheat.) 316, 423 (1819), alongside \textit{Bailey v.\ Drexel Furniture Co.\ (Child Labor Tax Case)}, 259 U.S.\ 20 (1922), Roberts held that the AAA was an unconstitutional scheme to regulate local agricultural production---a matter reserved to the states.

The Government contended that, unlike direct compulsion, the AAA's acreage-reduction program relied purely on voluntary contracts with farmers. Roberts offered two independent answers:
\begin{enumerate}
  \item \textbf{Economic Coercion:} Compliance was not genuinely voluntary. Refusing the subsidy left a farmer facing subsidized competitors who could undersell him, threatening ``financial ruin'' (297 U.S.\ at 71). Relying on \textit{Frost \& Frost Trucking Co.\ v.\ Railroad Comm'n}, 271 U.S.\ 583 (1926), Roberts declared that ``[t]his is coercion by economic pressure. The asserted power of choice is illusory'' (297 U.S.\ at 71).
  \item \textbf{No ``Purchased'' Regulation of Reserved State Subjects:} Even if farmers' participation were wholly voluntary, Congress ``cannot invade state jurisdiction to compel individual action; no more can it purchase such action'' (id.\ at 73). Roberts drew a sharp (though contested) line between a conditional appropriation specifying how federal money must be spent and a statute effective ``only upon assumption of a contractual obligation to submit to a regulation which otherwise could not be enforced'' (id.\ at 73). Accepting the Government's theory, he warned, would allow Congress to tax and subsidize its way into controlling wages, hours, manufacturing, and industrial location nationwide, obliterating the states and converting Congress into ``a parliament of the whole people, subject to no restrictions save such as are self-imposed'' (id.\ at 78).
\end{enumerate}

\section{Concurring and Dissenting Opinions}
\paragraph{Dissenting Opinion of Justice Stone (joined by Justices Brandeis and Cardozo).}
Justice Stone authored one of the most celebrated dissents in Supreme Court history, attacking the majority's methodology, empirical assumptions, and internal doctrinal logic:
\begin{itemize}
  \item \textbf{Judicial Self-Restraint:} Stone opened with a rebuke of judicial overreach, warning that courts are concerned only with legislative \textit{power}, not \textit{wisdom}, and that ``the only check upon our own exercise of power is our own sense of self-restraint'' (297 U.S.\ at 78--79).
  \item \textbf{Severability of Tax and Spend:} Unlike the regulatory penalty struck down in the \textit{Child Labor Tax Case}, the AAA's processing tax itself exerted no regulatory compulsion on cotton mills; any influence on agriculture arose solely from how the revenue was spent (id.\ at 80).
  \item \textbf{No Economic Coercion:} Stone dismantled Roberts's coercion claim as ``groundless speculation'' (id.\ at 83). ``Threat of loss, not hope of gain, is the essence of economic coercion,'' he wrote, pointing to Department of Agriculture data showing that 33\% of cotton farms in 1934 and 13\% in 1935 declined to participate---precisely why Congress later felt compelled to pass the mandatory Bankhead Cotton Act (id.\ at 81--83).
  \item \textbf{Internal Contradiction on Conditional Spending:} Most devastatingly, Stone exposed the contradiction between adopting Hamilton's broad spending power and forbidding conditions that ensure the funds achieve their national purpose. Drawing on \textit{McCulloch v.\ Maryland}, Stone argued that ``[t]he spending power of Congress is in addition to the legislative power, and not subordinate to it'' (id.\ at 85). Stripping Congress of the ability to attach conditions would lead to ``absurd consequences'': Congress could hand out seeds without requiring them to be planted, give disaster relief without requiring sanitation measures, or fund land-grant colleges under the Morrill Act without requiring agriculture to be taught (id.\ at 83--86). The true constitutional checks on spending, Stone concluded, are that the purpose be truly national, that it not coerce action reserved to the states, and ``the conscience and patriotism of Congress and the Executive'' (id.\ at 87).
\end{itemize}

\section{Key Quotations}
\begin{enumerate}
  \item \textbf{Formalist Judicial Review (Roberts, J., 297 U.S.\ at 62):}
  \begin{quote}
  ``When an act of Congress is appropriately challenged in the courts as not conforming to the constitutional mandate, the judicial branch of the Government has only one duty---to lay the article of the Constitution which is invoked beside the statute which is challenged and to decide whether the latter squares with the former.''
  \end{quote}

  \item \textbf{Adoption of the Hamiltonian Spending Power (Roberts, J., 297 U.S.\ at 66):}
  \begin{quote}
  ``While, therefore, the power to tax is not unlimited, its confines are set in the clause which confers it, and not in those of \S\ 8 which bestow and define the legislative powers of the Congress. It results that the power of Congress to authorize expenditure of public moneys for public purposes is not limited by the direct grants of legislative power found in the Constitution.''
  \end{quote}

  \item \textbf{Economic Coercion and the Ban on ``Purchased'' Compliance (Roberts, J., 297 U.S.\ at 71, 73):}
  \begin{quote}
  ``The power to confer or withhold unlimited benefits is the power to coerce or destroy. \dots\ This is coercion by economic pressure. The asserted power of choice is illusory. \dots\ The Congress cannot invade state jurisdiction to compel individual action; no more can it purchase such action.''
  \end{quote}

  \item \textbf{Judicial Self-Restraint (Stone, J., dissenting, 297 U.S.\ at 78--79):}
  \begin{quote}
  ``The power of courts to declare a statute unconstitutional is subject to two guiding principles of decision which ought never to be absent from judicial consciousness. One is that courts are concerned only with the power to enact statutes, not with their wisdom. The other is that, while unconstitutional exercise of power by the executive and legislative branches of the government is subject to judicial restraint, the only check upon our own exercise of power is our own sense of self-restraint. For the removal of unwise laws from the statute books appeal lies not to the courts, but to the ballot and to the processes of democratic government.''
  \end{quote}

  \item \textbf{Conditional Spending and Coercion (Stone, J., dissenting, 297 U.S.\ at 81, 85):}
  \begin{quote}
  ``Threat of loss, not hope of gain, is the essence of economic coercion. \dots\ It is a contradiction in terms to say that there is power to spend for the national welfare while rejecting any power to impose conditions reasonably adapted to the attainment of the end which alone would justify the expenditure.''
  \end{quote}
\end{enumerate}

\section{Doctrinal Significance and Modern Legacy}
\textit{United States v.\ Butler} occupies an ironic dual position in constitutional law: its narrow Tenth Amendment holding against the AAA was effectively gutted within seventeen months, yet its Hamiltonian construction of the General Welfare Clause---and its warning that financial inducement can cross into unconstitutional coercion---remains foundational through 2026.

\begin{itemize}
  \item \textbf{The 1937 Constitutional Revolution:} During the ``switch in time'' of 1937, the Supreme Court upheld the Social Security Act in \href{https://supreme.justia.com/cases/federal/us/301/548/}{\textit{Steward Machine Co.\ v.\ Davis}, 301 U.S.\ 548 (1937)} (unemployment compensation tax-credit scheme) and \href{https://supreme.justia.com/cases/federal/us/301/619/}{\textit{Helvering v.\ Davis}, 301 U.S.\ 619 (1937)} (old-age benefits). Justice Cardozo's opinions retained \textit{Butler}'s Hamiltonian reading of the Spending Clause while effectively adopting Stone's dissent on deference and non-coercion, holding that Congress's judgment of what constitutes the ``general welfare'' is controlling unless ``plainly arbitrary'' (\textit{Helvering}, 301 U.S.\ at 640) and distinguishing \textit{Butler} because Social Security taxes went into the general Treasury without earmarking.
  \item \textbf{The Modern Framework --- \textit{South Dakota v.\ Dole} (1987):} In \href{https://supreme.justia.com/cases/federal/us/483/203/}{\textit{South Dakota v.\ Dole}, 483 U.S.\ 203 (1987)} (7--2), Chief Justice Rehnquist expressly reaffirmed \textit{Butler}'s Hamiltonian premise that Congress may achieve objectives not thought to be within Article I's other enumerated fields through conditional grants, while synthesizing five limits on conditional spending: (1) pursuit of the general welfare (with substantial deference to Congress); (2) unambiguous notice to states; (3) germaneness/relatedness to the federal interest in particular national projects; (4) no independent constitutional bar (inducing states to engage in unconstitutional conduct); and (5) no coercion passing ``the point at which pressure turns into compulsion.'' Upholding a 5\% withholding of federal highway funds for states with a drinking age below 21, \textit{Dole} relegated \textit{Butler}'s anti-purchased-regulation rule to history while preserving its coercion principle in theory.
  \item \textbf{Revival of the Coercion Limit --- \textit{NFIB v.\ Sebelius} (2012):} Seventy-six years after \textit{Butler}, the Supreme Court for the first time since 1936 struck down a federal spending condition as unconstitutionally coercive in \href{https://supreme.justia.com/cases/federal/us/567/519/}{\textit{National Federation of Independent Business v.\ Sebelius}, 567 U.S.\ 519 (2012)}. Chief Justice John Roberts (joined by Justices Breyer and Kagan on the controlling Spending Clause opinion, with four dissenters agreeing on coercion) held that threatening states with the loss of \textit{all} existing Medicaid funding---often exceeding 10\% of a state's entire budget---if they declined the Affordable Care Act's Medicaid expansion was not ``relatively mild encouragement,'' but ``economic dragooning'' and ``a gun to the head.'' Echoing \textit{Butler}'s core federalism anxiety, \textit{NFIB} emphasized that the Spending Clause cannot be used to commandeer state legislatures or erase political accountability, while also underscoring \textit{Butler}'s distinction between a permissible tax (the ACA individual mandate, upheld as a tax) and an impermissible regulatory penalty.
  \item \textbf{Contract-Analogy Tightening Through 2025--2026:} In the mid-2020s, the Court has increasingly enforced \textit{Butler}'s structural federalism concerns not by reviving the Tenth Amendment ``reserved powers'' bar, but by rigorously policing the \textit{contract analogy} underlying Spending Clause legislation (see \href{https://www.everycrsreport.com/reports/R46827.html}{CRS Report R46827, updated Oct.\ 7, 2026}). In \href{https://supreme.justia.com/cases/federal/us/606/23-1275/}{\textit{Medina v.\ Planned Parenthood South Atlantic}, 606 U.S.\ 357 (2025)} (6--3), the Court narrowed private \S\ 1983 enforcement of Spending Clause conditions in the Medicaid Act, insisting on unmistakably clear rights-creating language before subjecting states to private suits. Even more pointedly, in \href{https://supreme.justia.com/cases/federal/us/609/23-1197/}{\textit{Landor v.\ Louisiana Department of Corrections and Public Safety}, 609 U.S.\ 59 (2026)} (6--3), Justice Gorsuch held that the Religious Land Use and Institutionalized Persons Act (RLUIPA)---enacted under the Spending Clause---cannot authorize personal-capacity money damages against individual state prison officials who did not personally, knowingly, and voluntarily consent to the federal funding bargain. Together, \textit{NFIB}, \textit{Medina}, and \textit{Landor} demonstrate that while Hamilton won the General Welfare Clause debate in \textit{Butler}, Justice Roberts's warning against an unbounded Spending Power swallowing state sovereignty remains a live, actively enforced doctrinal constraint in 2026.
\end{itemize}

\section{Pedagogical Critiques and Exam Hypotheticals}
\paragraph{Pedagogical Critiques.}
\begin{enumerate}
  \item \textbf{``Slot-Machine Jurisprudence'':} Generations of scholars have mocked Justice Roberts's assertion that judicial review merely requires laying the Constitution beside the statute to see whether they ``square.'' Because neither the General Welfare Clause nor the Tenth Amendment self-executing textually defines ``coercion'' or ``general welfare,'' Roberts's opinion illustrates how formalist rhetoric can mask deeply contested structural value judgments.
  \item \textbf{The Internal Tension of Half-Hearted Hamiltonianism:} As Justice Stone pointed out, adopting Hamilton's view that Congress may spend outside its other enumerated powers makes little sense if Congress is simultaneously barred from attaching conditions to ensure the money achieves its intended purpose. Roberts tried to preserve a distinction between ``conditional appropriations'' and ``contractual submission to regulation'' (297 U.S.\ at 73), but never explained why an ongoing funding condition is constitutionally different from a bilateral contract.
  \item \textbf{Earmarking Formalism and Standing:} Under \textit{Butler} and \textit{Steward Machine}, whether a taxpayer has standing---and whether a tax looks like a ``regulation''---turned heavily on whether Congress formally earmarked the tax revenue in the same statute rather than routing it through the general Treasury. Critics note that money is fungible, making the earmarking distinction a formalistic drafting trap rather than a substantive federalism safeguard.
\end{enumerate}

\paragraph{Exam Hypothetical 1 (Baseline --- Modern \textit{Butler}/\textit{Dole}/\textit{NFIB} Analysis).}
Suppose Congress enacts the ``National Soil Conservation and Carbon Reduction Act,'' imposing a 2\% federal excise tax on all commercial fertilizer sales (deposited into the general Treasury) and simultaneously establishing a \$15 billion annual grant program offering states federal agricultural-extension grants conditioned on each state enacting a statutory cap on nitrogen runoff from local farms. If a state declines to enact the cap, it forfeits 100\% of its existing federal Medicaid and highway funding (totaling 18\% of the state's annual budget).
\begin{itemize}
  \item \textbf{Analysis:} Under \textit{Butler}'s Hamiltonian holding and \textit{Helvering}, Congress may spend for national environmental and agricultural welfare even if local farming practices were beyond another enumerated power (and depositing the excise tax into the general Treasury avoids \textit{Butler}'s earmarking/standing pitfall). However, the funding condition fails both \textit{Dole} and \textit{NFIB}: (a) tying a fertilizer/nitrogen-runoff cap to \textit{Medicaid} funding fails \textit{Dole}'s relatedness requirement; and (b) threatening 18\% of a state's entire budget by clawing back massive, pre-existing programs is ``economic dragooning'' and ``a gun to the head'' under \textit{NFIB} (echoing \textit{Butler}'s coercion rationale).
\end{itemize}

\paragraph{Exam Hypothetical 2 (``Opposite Side'' Contrast --- Lawful Conditional Spending vs.\ Unlawful Third-Party Liability).}
Now modify Hypothetical 1 on the \textbf{opposite side} of the doctrinal line: Congress instead offers states a brand-new \$2 billion federal ``Clean Waterways Grant'' specifically to fund state agricultural-extension monitoring of nitrogen runoff. Participation is entirely optional; states that decline lose only the new \$2 billion grant (less than 0.3\% of any state's budget), while states that accept receive unambiguous statutory notice of the monitoring rules. However, the statute also purports to create a private federal cause of action for money damages against \textit{individual county farm-extension employees} in their personal capacities if they fail to enforce the runoff metrics.
\begin{itemize}
  \item \textbf{Analysis:} On the \textit{programmatic condition} itself, Congress wins cleanly under \textit{Butler}'s Hamiltonian holding, Stone's dissent, \textit{Dole}, and \textit{NFIB}: the condition is unambiguous, directly germane to the new grant's purpose, and involves no threat to existing funds (pure ``hope of gain,'' not ``threat of loss''). Yet on the \textit{remedial enforcement} side, the personal-capacity damages provision against individual state employees fails under \textit{Landor v.\ Louisiana Department of Corrections and Public Safety}, 609 U.S.\ 59 (2026): because Spending Clause legislation operates like a contract, Congress cannot bind non-consenting third-party individual employees to personal damages liability merely because their state employer accepted federal funds.
\end{itemize}

\end{document}
